\documentclass[a4paper,12pt]{article}
\usepackage[T2A]{fontenc}
\usepackage[utf8]{inputenc}
\usepackage[english,russian]{babel}
\usepackage{array}
\usepackage{booktabs}
\usepackage{siunitx}
\usepackage{tabularx}
\usepackage{multirow}
\usepackage{threeparttable}
\usepackage{longtable}  
\usepackage{float}
\usepackage{xcolor}
\usepackage{colortbl}

\title{Лабораторная работа №5 \\ Таблицы в LaTeX}
\author{Сунь Маосин}
\date{2026}

\begin{document}

\maketitle

\section{Цель работы}

Изучение возможностей LaTeX по созданию и форматированию таблиц.

\section{Простые таблицы}

\subsection{Базовая таблица с линиями}

\begin{tabular}{|l|c|r|}
  \hline
  Левый & Центр & Правый \\
  \hline
  1 & 2 & 3 \\
  \hline
  4 & 5 & 6 \\
  \hline
  7 & 8 & 9 \\
  \hline
\end{tabular}

\subsection{Таблица без вертикальных линий}

\begin{tabular}{lcr}
  \hline
  Левый & Центр & Правый \\
  \hline
  1 & 2 & 3 \\
  4 & 5 & 6 \\
  7 & 8 & 9 \\
  \hline
\end{tabular}

\section{Таблицы с пакетом booktabs}

Пакет \texttt{booktabs} создаёт профессионально выглядящие таблицы.

\begin{tabular}{lcr}
  \toprule
  Город & Население (млн) & Страна \\
  \midrule
  Москва & 12.5 & Россия \\
  Санкт-Петербург & 5.4 & Россия \\
  Новосибирск & 1.6 & Россия \\
  \bottomrule
\end{tabular}

\subsection{Разные виды линий booktabs}

\begin{tabular}{lcr}
  \toprule[2pt]
  Город & Население & Страна \\
  \midrule[1pt]
  Москва & 12.5 & Россия \\
  \cmidrule{1-2}
  Санкт-Петербург & 5.4 & Россия \\
  \addlinespace
  Новосибирск & 1.6 & Россия \\
  \bottomrule[2pt]
\end{tabular}

\section{Объединение ячеек}

\subsection{Горизонтальное объединение (multicolumn)}

\begin{tabular}{|l|c|r|}
  \hline
  \multicolumn{3}{|c|}{Заголовок на всю таблицу} \\
  \hline
  Левый & Центр & Правый \\
  \hline
  1 & 2 & 3 \\
  \hline
  \multicolumn{2}{|c|}{Объединение двух} & 3 \\
  \hline
\end{tabular}

\subsection{Вертикальное объединение (multirow)}

\begin{tabular}{|l|c|r|}
  \hline
  \multirow{2}{*}{Объединение} & A & 1 \\
  \cline{2-3}
  & B & 2 \\
  \hline
  C & D & 3 \\
  \hline
\end{tabular}

\subsection{Комбинация объединений}

\begin{tabular}{|l|c|c|}
  \hline
  \multicolumn{3}{|c|}{Заголовок} \\
  \hline
  \multirow{2}{*}{Группа} & \multicolumn{2}{c|}{Данные} \\
  \cline{2-3}
  & Значение 1 & Значение 2 \\
  \hline
  1 & 10 & 20 \\
  2 & 30 & 40 \\
  \hline
\end{tabular}

\section{Фиксированная ширина столбцов}

\subsection{Тип p (параграф)}

\begin{tabular}{|l|p{4cm}|}
  \hline
  Тип & Описание \\
  \hline
  p & Столбец с фиксированной шириной, текст переносится \\
  \hline
\end{tabular}

\subsection{Типы m и b (вертикальное выравнивание)}

\begin{tabular}{|l|m{3cm}|b{3cm}|}
  \hline
  Выравнивание & По центру (m) & По низу (b) \\
  \hline
  Текст & Этот текст выровнен по центру ячейки & Этот текст прижат к низу \\
  \hline
\end{tabular}

\subsection{Свои типы столбцов}

\begin{tabular}{|>{\bfseries}l|>{\itshape}c|>{\small}r|}
  \hline
  Жирный & Курсив & Мелкий \\
  \hline
  Текст & Текст & Текст \\
  \hline
\end{tabular}

\section{Числовое выравнивание с пакетом siunitx}

Пакет \texttt{siunitx} выравнивает числа по десятичной точке.

\begin{tabular}{S}
  {Число} \\
  123 \\
  45.67 \\
  8.9 \\
  0.1234 \\
\end{tabular}

\begin{tabular}{lS}
  \toprule
  Страна & {Население (млн)} \\
  \midrule
  Россия & 146.7 \\
  Китай & 1412.5 \\
  Индия & 1380.4 \\
  США & 331.9 \\
  \bottomrule
\end{tabular}

\section{Таблицы с фиксированной шириной (tabularx)}

Пакет \texttt{tabularx} позволяет задать общую ширину таблицы.

\begin{tabularx}{\textwidth}{|l|X|X|}
  \hline
  \multicolumn{1}{|c|}{Столбец 1} & \multicolumn{1}{c|}{Столбец 2} & \multicolumn{1}{c|}{Столбец 3} \\
  \hline
  Тип X & Этот столбец автоматически растягивается & Чтобы заполнить всю ширину \\
  \hline
  Ещё пример & Длинный текст, который будет перенесён & В столбце типа X \\
  \hline
\end{tabularx}

\subsection{Разная ширина столбцов}

\begin{tabularx}{\textwidth}{|l|X|X|X|}
  \hline
  Фикс. & Гибкий 1 & Гибкий 2 & Гибкий 3 \\
  \hline
  A & Длинный текст с переносом & Тоже гибкий & И этот тоже \\
  \hline
\end{tabularx}

\section{Многостраничные таблицы (longtable)}

Для длинных таблиц, которые не помещаются на одну страницу, используется пакет \texttt{longtable}.

\begin{longtable}{|l|c|c|}
  \caption{Длинная таблица, занимающая несколько страниц} \\
  \hline
  \textbf{№} & \textbf{Данные 1} & \textbf{Данные 2} \\
  \hline
  \endfirsthead
  
  \multicolumn{3}{c}{{\tablename\ \thetable{} -- продолжение}} \\
  \hline
  \textbf{№} & \textbf{Данные 1} & \textbf{Данные 2} \\
  \hline
  \endhead
  
  \hline
  \multicolumn{3}{r}{{Продолжение на следующей странице}} \\
  \endfoot
  
  \hline
  \endlastfoot
  
  1 & Значение A1 & Значение B1 \\
  2 & Значение A2 & Значение B2 \\
  3 & Значение A3 & Значение B3 \\
  4 & Значение A4 & Значение B4 \\
  5 & Значение A5 & Значение B5 \\
  6 & Значение A6 & Значение B6 \\
  7 & Значение A7 & Значение B7 \\
  8 & Значение A8 & Значение B8 \\
  9 & Значение A9 & Значение B9 \\
  10 & Значение A10 & Значение B10 \\
  11 & Значение A11 & Значение B11 \\
  12 & Значение A12 & Значение B12 \\
  13 & Значение A13 & Значение B13 \\
  14 & Значение A14 & Значение B14 \\
  15 & Значение A15 & Значение B15 \\
  16 & Значение A16 & Значение B16 \\
  17 & Значение A17 & Значение B17 \\
  18 & Значение A18 & Значение B18 \\
  19 & Значение A19 & Значение B19 \\
  20 & Значение A20 & Значение B20 \\
\end{longtable}

\section{Таблицы с примечаниями (threeparttable)}

Пакет \texttt{threeparttable} добавляет сноски к таблицам.

\begin{table}[H]
  \centering
  \begin{threeparttable}
    \caption{Таблица с примечаниями}
    \begin{tabular}{lS}
      \toprule
      Страна & {ВВП (трлн \$)\tnote{1}} \\
      \midrule
      США & 25.5 \\
      Китай & 18.3 \\
      Япония & 4.9 \\
      Германия & 4.3 \\
      \bottomrule
    \end{tabular}
    \begin{tablenotes}
      \item[1] Данные за 2023 год
      \item[2] Источник: Всемирный банк
    \end{tablenotes}
  \end{threeparttable}
\end{table}

\section{Цветные таблицы (опционально)}

С пакетом \texttt{xcolor} можно добавлять цвета.

\rowcolors{2}{gray!10}{white}

\begin{tabular}{lcr}
  \toprule
  Город & Население & Страна \\
  \midrule
  Москва & 12.5 & Россия \\
  Санкт-Петербург & 5.4 & Россия \\
  Новосибирск & 1.6 & Россия \\
  \bottomrule
\end{tabular}

\section{Вывод}

В ходе работы были изучены:

\begin{itemize}
  \item создание простых таблиц с помощью окружения \texttt{tabular};
  \item использование пакетов \texttt{array}, \texttt{booktabs}, \texttt{siunitx};
  \item объединение ячеек по горизонтали и вертикали;
  \item столбцы с фиксированной шириной;
  \item числовое выравнивание по десятичной точке;
  \item таблицы с заданной общей шириной;
  \item многостраничные таблицы;
  \item таблицы с примечаниями.
\end{itemize}

\end{document}